\chapter{Langages managés et fonctionnels sur les desks}\label{ch:ll:managed-and-functional-languages-on-the-desk}

En 2011, une bourse de détail a décrit publiquement comment sa logique d'appariement traitait six millions d'ordres par seconde sur un seul
fil d'exécution, écrite en Java (Fowler, 2011). Ses ingénieurs n'avaient pas échappé au \omterm{def:ll:managed-and-functional-languages-on-the-desk:gc}{ramasse-miettes}~; ils avaient conçu le système en fonction de lui~:
tout était pré-alloué, des tampons circulaires (ring buffers) de plusieurs millions d'emplacements créés au démarrage, aucun verrou dans la logique métier. Quelques
années plus tôt, une société de trading pour compte propre avait expliqué pourquoi elle écrivait ses systèmes de trading critiques en OCaml, un langage fonctionnel
doté de son propre \omterm{def:ll:managed-and-functional-languages-on-the-desk:gc}{ramasse-miettes} (Minsky et Weeks, 2008). Les langages managés sont présents sur les desks de trading, parfois sur
le chemin critique (hot path). Ce chapitre explique ce que font leurs environnements d'exécution et que ne font pas le C++ ni le Rust, comment un programme
sensible à la latence, écrit dans ces langages, est structuré, et ce que les sociétés qui les utilisent ont dit publiquement à ce sujet.

\section{Ramasse-miettes et ses pauses}

\begin{definition}[Ramasse-miettes, pause stop-the-world]\label{def:ll:managed-and-functional-languages-on-the-desk:gc}
Un \emph{ramasse-miettes}\index{ramasse-miettes} (garbage collector) récupère la mémoire que le programme ne peut plus atteindre, en parcourant
les références à partir des racines du programme (piles, variables globales). La plupart divisent le tas en générations~: les objets nouveaux vont dans une petite jeune génération,
collectée souvent et à faible coût~; les survivants sont promus dans une vieille génération, collectée rarement et à un coût qui croît
avec sa taille. Une \emph{pause stop-the-world}\index{pause stop-the-world} est un intervalle pendant lequel le ramasse-miettes
a arrêté tous les threads de l'application.
\end{definition}

Le compromis est explicite~: le programme ne libère jamais la mémoire et ne la libère jamais deux fois~; en contrepartie, son exécution
suit un rythme fixé par ses allocations. Les \omterm{def:ll:managed-and-functional-languages-on-the-desk:gc}{ramasse-miettes} modernes raccourcissent les pauses en effectuant l'essentiel de leur travail
en même temps que le programme~; un \omterm{def:ll:managed-and-functional-languages-on-the-desk:gc}{ramasse-miettes} Java très répandu visait d'abord des pauses inférieures à dix millisecondes, puis
plus tard, moins d'une milliseconde passée aux points où tous les threads s'arrêtent (encadré daté ci-dessous). Chacune de ces valeurs est énorme face
à un budget d'une microseconde, et le travail concurrent lui-même prend aux cœurs et à la bande passante mémoire du programme.

\begin{omfigure}
\begin{tikzpicture}[font=\footnotesize]
\draw[fill=omDef!10,draw=omDef,rounded corners=2pt] (0,0) rectangle (3.2,1.2);
\node[anchor=north] at (1.6,1.15) {jeune génération};
\node at (1.6,0.45) {nouveaux objets};
\draw[fill=omThm!10,draw=omThm,rounded corners=2pt] (5.0,0) rectangle (11.2,1.2);
\node[anchor=north] at (8.1,1.15) {vieille génération};
\node[align=center] at (8.1,0.42) {objets à longue durée de vie~:\\ le carnet d'ordres, les caches, l'état};
\draw[-Stealth,thick] (3.3,0.6) -- node[above]{survivants} node[below]{promus} (4.9,0.6);
\draw[-Stealth,thick] (-1.4,0.6) node[left]{allocation} -- (-0.1,0.6);
\draw[omDef,thick,decorate,decoration={brace,amplitude=4pt,mirror}] (0,-0.15) -- (3.2,-0.15) node[midway,below=5pt,align=center]{collectée souvent~:\\coût $\propto$ survivants};
\draw[omThm,thick,decorate,decoration={brace,amplitude=4pt,mirror}] (5.0,-0.15) -- (11.2,-0.15) node[midway,below=5pt,align=center]{collectée rarement~:\\coût croît avec le tas vivant};
\end{tikzpicture}

\omcaption{Un tas générationnel. L'allocation remplit la jeune génération, dont les collectes coûtent en proportion de ce qui
survit~; les survivants passent dans la vieille génération, dont les collectes complètes parcourent tout ce qui est vivant, et arrêtent ou ralentissent le
programme d'autant plus longtemps que le tas grandit.}\label{fig:ll:managed-and-functional-languages-on-the-desk:heap}
\end{omfigure}

\begin{dated}{2026-09}{Objectifs de pause d'un ramasse-miettes Java moderne}\label{dat:ll:managed-and-functional-languages-on-the-desk:jvm}
Le ZGC d'OpenJDK a été proposé (JEP 333) dans l'objectif que les durées de pause du \omterm{def:ll:managed-and-functional-languages-on-the-desk:gc}{ramasse-miettes} ne dépassent pas 10~millisecondes,
ses phases stop-the-world limitées au parcours des racines. Une proposition ultérieure (JEP 376) a déplacé le parcours des piles de threads
vers une phase concurrente, dans l'attente que moins d'une milliseconde soit consacrée aux safepoints du collecteur
sur des machines usuelles.
\end{dated}

\section{Le style managé à faible latence~: pré-allocation, mémoire hors tas, chauffe}

\begin{definition}[Mémoire hors tas (off-heap)]\label{def:ll:managed-and-functional-languages-on-the-desk:offheap}
La \emph{mémoire hors tas}\index{mémoire hors tas} (off-heap) désigne la mémoire qu'un programme managé utilise hors de son tas collecté (tampons d'octets
directs, fichiers mappés en mémoire, tableaux natifs), organisée par le programme et invisible au \omterm{def:ll:managed-and-functional-languages-on-the-desk:gc}{ramasse-miettes}~: il ne la parcourt jamais,
ne la déplace ni ne la libère jamais.
\end{definition}

Le style décrit par la bourse évoquée plus haut découle de l'économie du \omterm{def:ll:managed-and-functional-languages-on-the-desk:gc}{ramasse-miettes}. N'allouer rien par message après
le démarrage~: pré-allouer chaque objet dont le chemin critique a besoin et le réutiliser (les pools du chapitre~6)~; garder les données volumineuses et durables
(carnets d'ordres, tampons circulaires, journaux) hors du tas ou dans des tableaux plats de types primitifs, afin que les collectes complètes aient peu à parcourir~;
dimensionner la jeune génération pour que le peu de déchets qui subsiste la remplisse rarement~; et faire tourner un \omterm{def:ll:managed-and-functional-languages-on-the-desk:gc}{ramasse-miettes} discret. Ce qui reste
est un programme dont le \omterm{def:ll:managed-and-functional-languages-on-the-desk:gc}{ramasse-miettes} n'a rien à faire pendant la journée de trading. La limite du style tient à sa fragilité~:
un seul appel de bibliothèque qui alloue, un seul entier encapsulé (boxed), et le \omterm{def:ll:managed-and-functional-languages-on-the-desk:gc}{ramasse-miettes} revient sur le chemin critique, sans que rien dans le langage ne le signale~;
l'allocation doit donc être mesurée, comme au chapitre~6.

\begin{proposition}[Le budget de ramasse-miettes]\label{prop:ll:managed-and-functional-languages-on-the-desk:budget}
Un programme qui alloue $b$ octets par message, à $r$ messages par seconde, avec une jeune génération de $Y$ octets, collecte sa
jeune génération toutes les $Y/(rb)$ secondes. Il collecte au plus $k$ fois au cours d'une session de durée $T$ si et seulement si
$b \le kY/(rT)$.
\end{proposition}

\begin{proof}
L'allocation remplit la jeune génération à raison de $rb$ octets par seconde et chaque collecte la vide, donc les collectes ont lieu toutes les
$Y/(rb)$ secondes et sont au nombre de $rbT/Y$ dans la session~; $rbT/Y \le k$ est la condition énoncée.
\end{proof}

Avec une jeune génération de 2~GiB, 200\,000 messages par seconde et une session de six heures et demie, le budget pour une
collecte par session est inférieur à un demi-octet par message~: en pratique, zéro.

\section{Compilation à la volée et désoptimisation}

\begin{definition}[Compilation à la volée (JIT), période de chauffe (warm-up)]\label{def:ll:managed-and-functional-languages-on-the-desk:jit}
Dans la \emph{compilation à la volée}\index{compilation à la volée} (just-in-time compilation), un environnement d'exécution commence par interpréter ou compiler légèrement un programme,
mesure quel code est chaud et comment il se comporte, puis compile le code chaud en code machine optimisé à partir de ce profil.
La \emph{période de chauffe}\index{période de chauffe} (warm-up period) est le temps, ou le nombre d'exécutions,
qui précède le moment où le chemin critique s'exécute en code optimisé.
\end{definition}

\begin{definition}[Désoptimisation]\label{def:ll:managed-and-functional-languages-on-the-desk:deopt}
La \emph{désoptimisation}\index{désoptimisation} est le repli de l'environnement d'exécution du code optimisé vers un niveau plus lent lorsqu'une
hypothèse faite par l'optimiseur à partir du profil est mise en défaut (un site d'appel qui n'avait vu qu'un type en voit un autre, un embranchement
jamais pris est emprunté)~; le code peut être recompilé plus tard.
\end{definition}

Un compilateur à la volée est une \omterm{def:ll:cpp-for-latency-iii-the-compiler:pgo}{optimisation guidée par profil} (chapitre~8) réalisée pendant l'exécution, avec le même risque~: le
profil est ce que le programme a vu jusque-là. Un processus de trading qui démarre le matin est le plus lent sur ses premiers
messages, précisément au moment où l'ouverture est la plus chargée, et un chemin rare emprunté pour la première fois à un moment
volatil (un type de message inhabituel, une suspension de cotation) peut déclencher une \omterm{def:ll:managed-and-functional-languages-on-the-desk:deopt}{désoptimisation} à ce moment précis. Les remèdes usuels sont la chauffe, en rejouant du trafic synthétique ou
enregistré dans le code réel avant la séance (avec des ordres envoyés à un simulateur, jamais au marché), et la stabilité
des types et des embranchements du chemin critique. La compilation anticipée supprime la \omterm{def:ll:managed-and-functional-languages-on-the-desk:jit}{période de chauffe}, au prix du profil.

\section{Les sociétés à langage fonctionnel}

La société de trading pour compte propre évoquée en introduction a décrit son choix dans un article de revue~: OCaml comme langage de développement principal,
plus de vingt programmeurs OCaml et des centaines de milliers de lignes de code, utilisé pour les systèmes de trading critiques, la recherche,
les logiciels système et l'administration, et apprécié parce qu'il permettait à la société de «~produire rapidement du code lisible, correct et efficace~»
(Minsky et Weeks, 2008). L'argument ne portait pas sur la vitesse par message mais sur la correction par ingénieur~: un système de types expressif
qui encode des invariants (l'état d'un ordre comme type somme dont le compilateur vérifie que tous les cas sont traités), des données immuables par
défaut, et un code assez court pour être relu par les personnes qui s'en servent pour trader. Un essai ultérieur de la même société soutenait que le
prochain langage qu'un programmeur apprend devrait être fonctionnel (Minsky, 2011). Le \omterm{def:ll:managed-and-functional-languages-on-the-desk:gc}{ramasse-miettes} d'OCaml est lui aussi générationnel, avec une petite
tas mineur collecté souvent et un tas majeur collecté de façon incrémentale~; l'OCaml sensible à la latence, comme le
Java sensible à la latence, est écrit pour allouer peu sur le chemin critique.

\section{Choisir un langage pour un desk}

Le choix est moins binaire que ne le laissent entendre les débats. Plusieurs langages coexistent généralement dans une même société~: C++ ou Rust là où
se négocie la nanoseconde, un langage géré ou fonctionnel pour le gros de la logique dont
le \omterm{def:ll:where-latency-comes-from:budget}{budget de latence} se situe entre la microseconde et la milliseconde, et Python pour la recherche et les opérations (le livre~15, chapitre~8). Les
questions qui tranchent sont mesurables~: le \omterm{def:ll:where-latency-comes-from:budget}{budget de latence} au centile qui compte~; la discipline d'allocation que l'équipe
peut tenir et tester~; le coût des bogues que chaque langage empêche ou autorise~; et la taille et l'expérience de
l'équipe. En Python, langage de recherche de ce livre, les leçons de ce chapitre s'appliquent directement, comme le montre le tutoriel.

\section{Tutoriel~: le ramasse-miettes dans une boucle de messages Python}\label{tut:ll:managed-and-functional-languages-on-the-desk:tut}

\begin{tutorial}
\textbf{Objectif.} Observer le \omterm{def:ll:managed-and-functional-languages-on-the-desk:gc}{ramasse-miettes} de CPython interrompre une boucle de messages, puis le supprimer par pré-allocation. \textbf{État final~:}
\cref{fig:ll:managed-and-functional-languages-on-the-desk:gc}.
\begin{tutsteps}
\item \textbf{Deux carnets.} Le premier conserve chaque ordre vivant dans un dictionnaire Python, en ajoutant deux ordres pour chacun retiré, au fil du remplissage du carnet durant la matinée~; 
le second écrit les ordres dans un tableau d'enregistrements NumPy pré-alloué.
\omcode{code/low-latency/10-managed-and-functional-languages-on-the-desk/python/ll_gc.py}{19}{38}{Un carnet de dictionnaires et un carnet pré-alloué.}\label{lst:ll:managed-and-functional-languages-on-the-desk:books}
\item \textbf{Chronométrer le \omterm{def:ll:managed-and-functional-languages-on-the-desk:gc}{ramasse-miettes}.} \texttt{GcWatch} enregistre une fonction de rappel que CPython appelle avant et après chaque
collecte, et consigne chaque pause, par génération, dans \texttt{firm.lathist}.
\omcode{code/firm/gcwatch/firm_gcwatch.py}{27}{50}{Chronométrage de chaque collecte au moyen des rappels du module gc.}\label{lst:ll:managed-and-functional-languages-on-the-desk:watch}
\item \textbf{Mesurer la boucle} sur 200\,000 messages avec \texttt{python bench\_gc.py}~; la boucle ne stocke que des horodatages,
de sorte que seul le travail propre au carnet s'exécute entre deux d'entre eux.
\end{tutsteps}
\textbf{Pour aller plus loin.} Conserver le carnet de dictionnaires, mais appeler \texttt{gc.freeze()} après la \omterm{def:ll:managed-and-functional-languages-on-the-desk:jit}{période de chauffe} puis de nouveau
tous les 100\,000 messages, et observer ce qui advient des collectes complètes~; borner le carnet de dictionnaires à 1\,000 ordres vivants et
expliquer pourquoi le \omterm{def:ll:managed-and-functional-languages-on-the-desk:gc}{ramasse-miettes} ne se déclenche alors jamais.
\end{tutorial}

\begin{omfigure}
\begin{tikzpicture}
\begin{axis}[width=10cm,height=5.6cm,ybar,bar width=8pt,ymode=log,log origin y=infty,ymin=100,ymax=100000000,
  symbolic x coords={dict book,pooled book},xtick=data,xticklabels={carnet à dictionnaire,carnet à pool},enlarge x limits=0.45,
  ylabel={ns par message (échelle logarithmique)},tick label style={font=\small},label style={font=\small},grid=major,
  grid style={gray!25},legend columns=4,legend style={font=\footnotesize,at={(0.5,-0.2)},anchor=north,draw=none,
  fill=none,/tikz/every even column/.append style={column sep=6pt}},table/col sep=comma]
\addplot[fill=omDef!70,draw=omDef] table[x=design,y=p50]{figdata/low-latency/10-managed-and-functional-languages-on-the-desk/measured_gc.csv};
\addplot[fill=omDef!35,draw=omDef] table[x=design,y=p99]{figdata/low-latency/10-managed-and-functional-languages-on-the-desk/measured_gc.csv};
\addplot[fill=omThm!45,draw=omThm] table[x=design,y=p999]{figdata/low-latency/10-managed-and-functional-languages-on-the-desk/measured_gc.csv};
\addplot[fill=gray!40,draw=gray] table[x=design,y=max]{figdata/low-latency/10-managed-and-functional-languages-on-the-desk/measured_gc.csv};
\legend{p50,p99,p99.9,maximum}
\end{axis}
\end{tikzpicture}

\omcaption{Une boucle de messages Python sur 200\,000 ordres~: un carnet de dictionnaires, qui a déclenché 259 collectes jeunes, 23 intermédiaires
et 2 complètes, et un carnet pré-alloué en NumPy, qui n'en a déclenché aucune. Le maximum du carnet de dictionnaires correspond à sa
collecte complète. CPython~3.10. Mesuré sur un portable (Intel Core Ultra 7 155H) sous WSL2, sans cœurs isolés. Données~:
\texttt{bench\_gc.py}.}\label{fig:ll:managed-and-functional-languages-on-the-desk:gc}
\end{omfigure}

Le résultat comporte deux volets, et tous deux comptent. Le carnet pré-alloué a supprimé le \omterm{def:ll:managed-and-functional-languages-on-the-desk:gc}{ramasse-miettes}~: aucune collecte sur 200\,000
messages, et une queue plus de vingt fois plus courte à son maximum. Il n'est pas plus rapide à la médiane~: écrire trois champs
d'un tableau NumPy depuis Python coûte plus que créer un petit dictionnaire. Et le pire message du carnet de dictionnaires a duré
aussi longtemps que sa dernière collecte complète, une vingtaine de millisecondes, qui a parcouru chaque ordre encore vivant. CPython ajoute
une subtilité qui lui est propre~: son \omterm{def:ll:managed-and-functional-languages-on-the-desk:gc}{ramasse-miettes} compte les allocations moins les désallocations, de sorte qu'une boucle dont l'ensemble vivant reste borné
ne déclenche jamais de collecte~; un carnet qui grossit, lui, en déclenche.

\section{Construction~: surveillance du ramasse-miettes et pool de messages}\label{bld:ll:managed-and-functional-languages-on-the-desk:gcwatch}

\begin{build}
\bfield{Rôle} Les services Python de la société (serveurs de recherche, reporting, pilotes du simulateur du livre~10)
mesurent leur \omterm{def:ll:managed-and-functional-languages-on-the-desk:gc}{ramasse-miettes} et le tiennent à l'écart de leurs boucles de messages.
\bfield{Interface} Python \texttt{firm\_gcwatch}~: \texttt{GcWatch()} utilisé comme gestionnaire de contexte avec \texttt{pauses[gen]}
(\texttt{LatHist}, en nanosecondes), \texttt{collections[gen]}, \texttt{total\_ns}~; \texttt{RecordPool(n, dtype)} avec
\texttt{next()} et \texttt{buf}~; \texttt{MESSAGE} (horodatage, prix, quantité, côté)~;
\texttt{steady\_state(step, n, warm)}.
\bfield{Règles} La fonction de rappel est retirée à la fin du contexte, même en cas d'erreur~; \texttt{steady\_state} gèle le tas après
la \omterm{def:ll:managed-and-functional-languages-on-the-desk:jit}{période de chauffe} et le dégèle ensuite~; le pool ne grandit jamais.
\bfield{Tests d'acceptation} \texttt{code/firm/gcwatch/tests/}~: une boucle à pool de 50\,000 messages ne provoque aucune collecte~; un
renouvellement borné n'en provoque pas non plus~; un carnet croissant provoque des collectes jeunes, chacune chronométrée.
\bfield{Pour aller plus loin} Mesurer le débit d'allocation (octets par message) avec \texttt{tracemalloc} dans une version de test~; déclencher une alerte si une
collecte complète survient pendant les heures de marché.
\end{build}

\begin{omsources}
\item M.~Fowler, ``The LMAX architecture'', 12 juillet 2011.
\item Y.~Minsky et S.~Weeks, ``Caml trading: experiences with functional programming on Wall Street'', \emph{Journal of
Functional Programming} 18(4), 2008~; Y.~Minsky, ``OCaml for the masses'', \emph{ACM Queue}, 2011.
\item OpenJDK, JEP 333 (ZGC) et JEP 376 (concurrent thread-stack processing).
\item Documentation Python, module \texttt{gc}.
\end{omsources}

\section{Exercices}

\begin{exercise}[$\star$]\label{exo:ll:managed-and-functional-languages-on-the-desk:1}
Un moteur Java alloue 64 octets par message à raison de 100\,000 messages par seconde, avec une jeune génération de 1~GiB. À quelle fréquence
collecte-t-il la jeune génération, et combien de fois le fait-il sur une session de 6.5 heures~?
\end{exercise}

\begin{exercise}[$\star$]\label{exo:ll:managed-and-functional-languages-on-the-desk:2}
La pause d'un collecteur dure \qty{1}{\milli\second}. Combien de messages arrivant à 200\,000 par seconde font la queue derrière elle, et combien de temps
attend le dernier d'entre eux~?
\end{exercise}

\begin{exercise}[$\star$]\label{exo:ll:managed-and-functional-languages-on-the-desk:3}
D'après la \cref{fig:ll:managed-and-functional-languages-on-the-desk:gc}, de combien le carnet à pool réduit-il le maximum et
le p99.9, et comment sa médiane se compare-t-elle~?
\end{exercise}

\begin{exercise}[$\star\star$]\label{exo:ll:managed-and-functional-languages-on-the-desk:4}
Pourquoi une boucle Python qui conserve les 1\,000 derniers ordres dans un deque borné ne déclenche jamais le \omterm{def:ll:managed-and-functional-languages-on-the-desk:gc}{ramasse-miettes}, alors qu'un carnet qui grossit
le déclenche~?
\end{exercise}

\begin{exercise}[$\star\star$]\label{exo:ll:managed-and-functional-languages-on-the-desk:5}
Une stratégie Java est la plus rapide après dix minutes de négociation et la plus lente à l'ouverture. Expliquer ce phénomène et proposer deux remèdes.
\end{exercise}

\begin{exercise}[$\star\star$]\label{exo:ll:managed-and-functional-languages-on-the-desk:6}
La bourse d'ouverture conservait des tampons circulaires de 20~millions d'emplacements côté entrée. À 6~millions d'ordres par seconde, combien de
secondes de flux d'entrée un tel tampon contient-il~?
\end{exercise}

\begin{exercise}[$\star\star\star$]\label{exo:ll:managed-and-functional-languages-on-the-desk:7}
\emph{Programmation.} Appeler \texttt{gc.freeze()} dans le carnet à dictionnaire après la \omterm{def:ll:managed-and-functional-languages-on-the-desk:jit}{période de chauffe} (warm-up), puis tous les 100\,000 messages. Compter les
collectes complètes et noter le message le plus lent. Que sacrifie le gel~?
\end{exercise}

\begin{exercise}[$\star\star\star$]\label{exo:ll:managed-and-functional-languages-on-the-desk:8}
\emph{Trouver la faille.} «~La passerelle Java utilise le \omterm{def:ll:managed-and-functional-languages-on-the-desk:gc}{ramasse-miettes} concurrent, avec des pauses inférieures à une milliseconde, donc le \omterm{def:ll:managed-and-functional-languages-on-the-desk:gc}{ramasse-miettes}
n'affecte pas notre latence.~»
\end{exercise}

\section{Problème~: Le budget du ramasse-miettes}

\begin{problem}[{Problème du week-end --- un moteur JVM un jour de forte activité}]\label{pb:ll:managed-and-functional-languages-on-the-desk:1}
Une société de trading fait tourner un moteur de routage d'ordres sur la JVM, avec une jeune génération de 2~GiB. Un jour ordinaire, il traite 200\,000 messages par
seconde pendant une session de 6.5~heures et alloue en moyenne 48~octets par message. Une revue de code montre que l'essentiel vient de
des entiers encapsulés (boxed integers) et d'un appel de journalisation qui met en forme une chaîne.

\medskip\noindent\textbf{Partie I --- Le budget.}
\begin{enumerate}
\item Combien d'octets le moteur alloue-t-il par seconde, et par session~?
\item À quelle fréquence collecte-t-il la jeune génération, et combien de fois par session~?
\item Quelle est l'allocation par message qui permet au plus une collecte de la jeune génération par session~?
\item Que signifie ce budget en pratique~?
\end{enumerate}

\medskip\noindent\textbf{Partie II --- Les pauses.}
\begin{enumerate}[resume]
\item Si chaque collecte de la jeune génération met le moteur en pause pendant \qty{0.5}{\milli\second}, combien de temps par session le moteur reste-t-il en pause~?
\item Combien de messages s'accumulent derrière une telle pause, et combien de temps attend celui qui se trouve en queue de file~?
\item Un jour où le trafic est trois fois plus important, comment les réponses changent-elles~?
\item Pourquoi le jour chargé est-il celui qui compte~?
\end{enumerate}

\medskip\noindent\textbf{Partie III --- Les correctifs.}
\begin{enumerate}[resume]
\item Par quoi remplacer les entiers encapsulés sur le chemin critique (hot path)~?
\item Par quoi remplacer la journalisation formatée (chapitre~23)~?
\item Où placer le carnet d'ordres et les tampons circulaires, et pourquoi~?
\item Comment vérifier, dans la suite de tests, que l'allocation par message est nulle~?
\end{enumerate}

\medskip\noindent\textbf{Partie IV --- Le verdict.}
\begin{enumerate}[resume]
\item Énoncer le \emph{résultat nommé}~: l'allocation par message qui limite le moteur à une collecte de la jeune génération par session,
et le temps de pause par session pour les 48~octets actuels.
\item Que reste-t-il du coût du \omterm{def:ll:managed-and-functional-languages-on-the-desk:gc}{ramasse-miettes} une fois l'allocation nulle~?
\item Pourquoi le moteur doit-il encore passer par sa \omterm{def:ll:managed-and-functional-languages-on-the-desk:jit}{période de chauffe} avant l'ouverture~?
\item Que montre l'expérience Python du tutoriel qui vaut aussi pour la JVM~?
\item Pourquoi la société utilisant un langage fonctionnel accordait-elle de la valeur à ce langage, selon son propre article~?
\item Dans quels cas faudrait-il choisir C++ ou Rust plutôt qu'un langage managé pour ce moteur~?
\item Que doit mesurer la société avant de trancher~?
\item En une phrase~: quel coût un \omterm{def:ll:managed-and-functional-languages-on-the-desk:gc}{ramasse-miettes} fait-il peser sur un chemin critique~?
\end{enumerate}
\end{problem}

\section{Questions d'entretien}

\begin{interviewq}[$\star$ \iqroles{développeur}]\label{iq:ll:managed-and-functional-languages-on-the-desk:1}
Comment fonctionne un \omterm{def:ll:managed-and-functional-languages-on-the-desk:gc}{ramasse-miettes} générationnel, et pourquoi est-il générationnel~?
\end{interviewq}

\begin{interviewq}[$\star\star$ \iqroles{développeur}]\label{iq:ll:managed-and-functional-languages-on-the-desk:2}
Comment écririez-vous un service à faible latence en Java pour que le \omterm{def:ll:managed-and-functional-languages-on-the-desk:gc}{ramasse-miettes} n'interfère pas~?
\end{interviewq}

\begin{interviewq}[$\star\star$ \iqroles{développeur}]\label{iq:ll:managed-and-functional-languages-on-the-desk:3}
Qu'est-ce que la \omterm{def:ll:managed-and-functional-languages-on-the-desk:jit}{période de chauffe} (warm-up) dans un environnement d'exécution compilé à la volée (JIT), et comment la gérez-vous avant l'ouverture~?
\end{interviewq}

\begin{interviewq}[$\star\star$ \iqroles{développeur}]\label{iq:ll:managed-and-functional-languages-on-the-desk:4}
Qu'est-ce que la \omterm{def:ll:managed-and-functional-languages-on-the-desk:deopt}{désoptimisation}, et pourquoi peut-elle survenir au pire moment~?
\end{interviewq}

\begin{interviewq}[$\star\star$ \iqroles{développeur, chercheur}]\label{iq:ll:managed-and-functional-languages-on-the-desk:5}
Pourquoi une société de trading choisirait-elle un langage fonctionnel pour ses systèmes de trading~?
\end{interviewq}

\begin{interviewq}[$\star\star\star$ \iqroles{développeur}]\label{iq:ll:managed-and-functional-languages-on-the-desk:6}
Le p99.9 d'un service Python est correct, mais son maximum atteint 20 millisecondes une fois toutes les quelques minutes. Comment savez-vous si
c'est le \omterm{def:ll:managed-and-functional-languages-on-the-desk:gc}{ramasse-miettes}, et que faites-vous à ce sujet~?
\end{interviewq}
